\section{Conditional Scoring and Context Purification}\label{sec:method}
\subsection{Task and information boundary}
Let $\widetilde X$ be an observed table with $n$ rows and $d$ eligible columns.
The detector outputs $S\in\mathbb R^{n\times d}$, with larger values indicating
greater suspicion. Evaluation compares observed strings with aligned clean strings;
this defines benchmark disagreement, rather than a universal semantic definition
of an error. Clean values and the error matrix are unavailable to deployment arms.
Only a separately labelled oracle receives clean contexts.

For seed $s$, partition rows into $K=5$ outer folds $Q_f$. Sample a sorted base
context $B_f$ of at most $C=512$ rows from $\{1,\ldots,n\}\setminus Q_f$.
Share these bases across target columns, removal arms and models. To score column
$j$, drop $j$ from the features and predict its observed values from the remaining
eligible columns. Excluded identifier columns are not conditioning features.
The query fold never enters a fitted deployment context. Other rows of the same
entity may be used in the primary row-level analysis.

Typing and feature encoding are fitted on the complete observed table for the
final pass. Numerical features use parsed numbers with unparseable entries as
missing; other features use observed frequency-rank codes and an unknown code.
This is label-free transductive preprocessing: outer queries can affect encoding,
although they do not enter the model's fitted context. We distinguish this
assumption from a fully inductive fit on context rows only.

\subsection{Mixed-type suspicion scores}
A column is numerical if at least $80\%$ of its nonempty observed strings parse
as finite numbers. Otherwise it is categorical, unless its observed unique-value
ratio exceeds $0.5$, in which case the target is unsupported. Unsupported cells
receive score zero and remain in every whole-universe evaluation denominator.
They may still supply encoded conditioning features.

For a parseable numerical query $y$ and context of size $b$, let $b_p$ be the
number of parseable context targets. With predictive CDF $F$, the score is
\begin{equation}
 S(y)=-\log\frac{b_p+0.5}{b+1}
       -\log\max\{2\min(F(y),1-F(y)),10^{-12}\}.
\end{equation}
The second term is a two-sided tail score, not a probability density. With fewer
than ten parseable context targets or no numerical range, a smoothed exact-match
frequency replaces that term. Unparseable query tokens use a smoothed
nonparseability probability followed by their smoothed frequency among
unparseable context targets. These rules retain empty strings and formatting
variants rather than repairing them before detection.

For a categorical target with at most ten context classes, the model predicts
class probabilities. With more classes, retain the nine most frequent and pool
the others into an OTHER class. Divide OTHER mass by smoothed within-pool
frequencies, allowing probability for an unseen query token. If the context
contains a single class, use the degenerate class distribution. A final
additive smoothing term assigns mass to values absent from the context:
\begin{equation}
 S(y)=-\log\left[\frac{b\widehat p(y\mid x_{-j})+1/(v+1)}{b+1}\right],
\end{equation}
where $v$ is the number of distinct observed context target values and
$\widehat p$ includes OTHER redistribution when applicable. Previously unseen
tokens share an unknown bucket; this is not a probability model over every
possible unseen string. These numerical
and categorical scores share an orientation but have no proven common
calibration. We therefore retain type-specific analyses alongside pooled AP.

\paragraph{Relation to residual and conformal scores.}
DDC standardizes residuals from robust linear relations among numerical
variables~\cite{rousseeuw2018ddc}; our scores instead come from the predictive
distribution of a pretrained or task-fitted model, without a linearity assumption.
CDC calibrates prediction sets on clean data~\cite{jager2024cdc}. Our deployment arms
have no clean calibration split and the context may contain errors, so $S$ is used
only as an uncalibrated ranking evaluated by precision--recall metrics. Removal is
based on first-pass suspicion scores, not on a gradient-based influence estimate.

\subsection{Outer-isolated selection}
The purification parameter is fixed at $q=0.15$. For each outer fold, score
every row of its training pool by an inner raw cross-fit. Inner typing,
encoding, contexts and fits all use that pool exclusively. Let $U^{(f)}_{ij}$
be the inner out-of-fold score for target $j$. For each base $B_f$, remove
the $r_f=\operatorname{round}(q|B_f|)$ rows with largest $U^{(f)}_{ij}$,
breaking ties by increasing row index. Fit the final conditional model on
the remaining dirty context and score $Q_f$. Selection is column specific:
the rows removed for one target need not be removed for another.

\begin{figure}[t]
\centering
\fbox{\begin{minipage}{0.91\linewidth}\small
\textbf{Outer-isolated purification for seed $s$}\par
1. Construct row folds $Q_f$ and shared bases $B_f$.\par
2. For each outer fold $f$, form $P_f=\{1,\ldots,n\}\setminus Q_f$.\par
3. Inner-cross-fit raw scores $U^{(f)}$ using only $\widetilde X[P_f]$.\par
4. For each target $j$, rank $B_f$ by $U^{(f)}_{ij}$ and remove $r_f$ rows.\par
5. Fit on the retained dirty context, omitting feature $j$; score $Q_f$.\par
6. Retain unsupported cells and failure masks in the full output universe.
\end{minipage}}
\caption{Selection has no dependency on the outer query through the first pass.
Final encoding remains explicitly transductive. This is one trimming step,
not an iterative procedure or a gradient-based influence estimate.}
\label{fig:algorithm}
\end{figure}

Reusing a raw score matrix computed on the entire table would break this
isolation: training examples used to score a prospective context row could
include the eventual outer query. The nested construction prevents that
dependency. It does not prove that the highest scores identify corrupted
contexts, nor that removing them preserves rare correct values.

\subsection{Controls, diagnostics and compute}
\emph{Raw} retains $B_f$. \emph{Random} removes $r_f$ rows without replacement
using a separate deterministic stream for each seed, fold and target.
\emph{Oracle} keeps the same base IDs but substitutes clean target and feature
values in the context; query values remain observed and the encoder remains
dirty fitted. Its gap is diagnostic and need not be positive.
\emph{Leaky}, on real-error conditions only, unions the query fold with the base;
its intentional violation of exclusion is reported separately.

The secondary \emph{raw-LEO} arm constructs folds of observed entity keys,
using provider number, flight and brewery identifier on their declared tables.
No observed query key enters a fitted context. Corrupted keys can split a true
entity and shared or empty keys can merge entities. Bases differ from row-level
bases, and the final encoder is still transductive; this arm is a boundary
sensitivity analysis rather than proof of independence.

Purification requires up to $K^2d$ additional column tasks for nested scoring,
besides up to $Kd$ final tasks. The corresponding number for a single raw or
random arm is at most $Kd$; degenerate or unsupported targets can bypass a model
fit, while LightGBM residual calibration can require additional fits.
These are task-count bounds, not measured neural forward-call counts.
We measure nested preparation and inclusive purification time separately.
The accounted matched-compute random cost is ordinary random time plus the
same nested preparation time, using unchanged random scores; it is not a separately
timed arm. No task-specific TabPFN weight
update is performed; its context fitting and preprocessing remain chargeable
work. LightGBM, in contrast, fits task-specific trees for each context.
